%\usepackage{xr}

% \externaldocument{note01}
% \externaldocument{note02}
% \externaldocument{note03}
% \externaldocument{note04}
% \externaldocument{note05}
% \externaldocument{note06}
% \externaldocument{note07}
% \externaldocument{note08}
% \externaldocument{note09}
% \externaldocument{note10}
% \externaldocument{note11}

\graphicspath{{./input/graphs/note07/}}


% \zexternaldocument{note01}
% \zexternaldocument{note02}
% \zexternaldocument{note03}
% \zexternaldocument{note04}
% \zexternaldocument{note05}
% \zexternaldocument{note06}
% \zexternaldocument{note07}
% \zexternaldocument{note08}
% \zexternaldocument{note09}
% \zexternaldocument{note10}
% \zexternaldocument{note11}
% \externaldocument{note12}



\week{8}
% \section*{Week 8. Lecture 19.}
\lecture{19}{Logarithm in $\mathbb C$}

\section{Function Ln $z$}

\subsection{Logarithm branches and properties}

\begin{definition}\label{definition:15.1}
     The natural logarithm of a complex number $z$ is a complex number $w$ such that $e^w = z$. Notation: $w = \text{Ln } z$.
        The logarithm of $r \in \mathbb{R}$ denote as usual $\ln r$.
\end{definition}

Let us develop a formula for the computation of the complex logarithm. Write: $z = r(\cos \varphi + i\sin \varphi), w = u + iv$.

By Def~\ref{definition:15.1}: $e^{u+iv} = r(\cos \varphi + i\sin \varphi) \implies e^u (\cos v + i\sin v) = r(\cos \varphi + i\sin \varphi)$.

By comparing the complex numbers, we conclude:
\begin{align*}
    &e^u = r \implies u = \ln r \\
    &\begin{cases}
    \cos v = \cos \varphi \\
    \sin v = \sin \varphi
    \end{cases} \implies v = \varphi + 2\pi k, k \in \mathbb{Z}
\end{align*}
$\implies \text{Ln } z = w = u + iv = \ln r + i(\varphi + 2\pi k)$ for $k \in \mathbb{Z}$.

Recall: $r = |z|, \varphi + 2\pi k = \arg z$, thus 
%(*)
\begin{equation*}\label{equation:ln_z_formula}
    \text{Ln } z = \ln |z| + i \arg z. \tag{*}
\end{equation*}

$\implies \text{Ln } z$ is defined for any $z \neq 0$ and it is multi-valued function (for different values of $k$ the function attains different values).

\begin{example}\label{example:17.1} 
    Compute:
    \begin{enumerate}
        \item $\text{Ln } i$
        \item $\text{Ln } (-4)$
    \end{enumerate}

    \begin{soln}
        \begin{enumerate}
            \item $\text{Ln } i = \ln |i| + i(\arg i + 2\pi k) = \ln 1 + i(\frac{\pi}{2} + 2\pi k) = i(\frac{\pi}{2} + 2\pi k)$.
            \item $\text{Ln } (-4) = \ln |-4| + i(\arg(-4) + 2\pi k) = \ln 4 + i(\pi + 2\pi k)$.
        \end{enumerate}
    \end{soln}
\end{example}
\begin{example} \label{example:17.2}
     Find an analytic function $f(z)$ such that $\text{Re } f(z) = \arctan \frac{y}{x}$ and $f(-i) = \frac{3\pi}{2} + 2i$.
\begin{soln}
    
    In Ex~\ref{example:11.3} we have seen that $f(z) = \arctan \frac{y}{x} - \frac{i}{2} \ln(x^2 + y^2) + ic$.
    
    $\implies$ the only question is to determine the value of the constant $c$.
    
    Let $z = x + iy$. Recall: $\arg z - \theta$ s.t. $\cos \theta = \frac{x}{|z|}, \sin \theta = \frac{y}{|z|} \implies \frac{\sin \theta}{\cos \theta} = \tan \theta = \frac{y}{x} \implies \theta = \arctan \frac{y}{x} \implies \arg z = \arctan \frac{y}{x}$ (up to $2\pi$), $|z|^2 = x^2 + y^2$.
    $\implies f(z) = \arg z - \frac{i}{2} \ln |z|^2 + ic = -i(\ln|z| + i \arg z) + ic$, or $f(z) = -i \text{Ln } z + c_1$, where $c_1 \in \mathbb{C}$ is a constant, let us compute it: $\frac{3\pi}{2} + 2i = f(-i) = -i\text{Ln }(-i) + c_1 = -i(0 + i\frac{3\pi}{2}) + c_1 = \frac{3\pi}{2} + c_1 \implies c_1 = 2i$ and $f(z) = i\text{Ln } z + 2i$.
\end{soln}
\end{example}

\begin{proposition}\label{proposition:15.1} For any $z_1, z_2 \in \mathbb{C}, n \in \mathbb{N}$:
\begin{enumerate}
    \item $\text{Ln}(z_1 z_2) = \text{Ln } z_1 + \text{Ln } z_2$
    \item $\text{Ln}(z_1/z_2) = \text{Ln } z_1 - \text{Ln } z_2$
    \item $\text{Ln } z^n = n \text{Ln } z$
    \item $\text{Ln} \sqrt[n]{z} = \frac{1}{n} \text{Ln } z$
\end{enumerate}
\begin{proof}
By formula (\ref{equation:ln_z_formula}), Prop~\ref{proposition:1.2}, and Prop~\ref{proposition:1.3}, we get:
$\text{Ln}(z_1 z_2) = \ln|z_1 z_2| + i\arg(z_1 z_2) = \ln(|z_1||z_2|) + i(\arg z_1 + \arg z_2) = (\ln|z_1| + i\arg z_1) + (\ln|z_2| + i\arg z_2) = \text{Ln } z_1 + \text{Ln } z_2$.
The rest in the same way. (FINISH!)
\end{proof}
\end{proposition}
\begin{remark}
    If $z_1 = r_1 e^{i\varphi_1}, z_2 = r_2 e^{i\varphi_2}$, then 1) can be written as $\text{Ln}(z_1 z_2) = \ln r_1 + \ln r_2 + i\varphi$, where $\varphi = \varphi_1 + \varphi_2 \pmod{2\pi}$.
\end{remark}

Q-n: \checkthis{}- 2), 3), 4)!

Now let us study $\text{Ln } z$ in the same spirit as we studied the complex root.

$w = \text{Ln } z$: For any $z$ correspond infinite number of values $w$. We choose a single-valued branch by fixing $k_0$: $w_{k_0} = \ln r + i(\varphi + 2\pi k_0)$.

\includegraphics{img_graph1.pdf}

Let $z$ move on some closed curve that does not encircle the origin, $z = 0$. After completing a full circle, $z$ returns to the same point with the same values of $r$ and $\varphi \implies$ the point $w_{k_0}$ does not change.

\includegraphics{img_graph2.pdf}

If $z$ moves (anticlockwise) on some closed curve that encircles the origin, then when $z$ returns to the starting point the modulus $r$ is the same, but the argument of $z$ will increase by $2\pi$ and $\text{Ln } z$ will attain a new value: $w_{new} = \ln r + i((\varphi + 2\pi k_0) + 2\pi) = \ln r + i(\varphi + 2\pi(k_0 + 1)) = w_{k_0 + 1}$.

So, after completing a full circle around the origin, the branch $w_{k_0}$ "jumps" to a new branch $w_{k_0 + 1}$. If two full circles are completed: $\text{Ln } z$ will attain another (new) value: $w_{k_0 + 2}$, since the modulus is the same, but the argument will increase by $4\pi$. In the same way, by completing finite number of full circles (anticlockwise) around the origin it is possible to pass from $w_{k_0}$ branch of $\text{Ln } z$ to any given in advance branch. For example, after 7 full circles around the origin $w_{k_0}$ will pass to $w_{k_0 + 7}$ (if in anticlockwise direction) or to $w_{k_0 - 7}$ (if clockwise direction).

The point $z = 0$ is the branch point of the function $\text{Ln } z$.

Conclusion: In any domain that does not contain closed curves
that encircle $z = 0$ one can choose infinitely many single-valued continuous branches of multi-valued function $w = \text{Ln } z$. If we cut the complex plane along any ray that starts at $z = 0$ (including $z = 0$) - for example, the negative half of the real axis including $z = 0$, we will get a cut-plane in which one can choose single-valued continuous branches. 
\includegraphics{img_graph3.pdf}

In this case the branch $w_{k_0}$ maps the cut-plane $z$ to a strip $\{w: k_0 \pi < \text{Im } w < (k_0 + 2)\pi\}$ in the w-plane.

\checkthis{} $z = z_0$ is the branch point of $\text{Ln}(z - z_0)$ and in the cut-plane with the cut along some ray starting at $z = z_0$ (including $z_0$) one can choose infinitely many single-valued continuous branches.

\begin{proposition}\label{proposition:15.2} Every branch of the function $\text{Ln } z$ is an analytic function in the cut-plane (cutted along some ray starting at the origin) and $(\text{Ln } z)' = \frac{1}{z}$.

    \begin{proof}
We prove that $\text{Ln } z$ is differentiable. Let $\Delta z$ be such that $z$ and $z + \Delta z$ are in the same domain. We compute:
\begin{align*}
 &\frac{\text{Ln}(z + \Delta z) - \text{Ln } z}{\Delta z} = \frac{\Delta w}{\Delta z} = \frac{\Delta w}{e^{w + \Delta w} - e^w} = \frac{1}{\frac{e^{w + \Delta w} - e^w}{\Delta w}} \\
    \Delta z &= e^{w + \Delta w} - e^w \\
    \text{Ln } z &= \text{Ln } e^w = w \implies \text{Ln}(z + \Delta z) = \text{Ln}(e^w e^{\Delta w}) = w + \Delta w \text{ and } \\
    &\implies \text{Ln}(z + \Delta z) - \text{Ln } z = w + \Delta w - w = \Delta w
\end{align*}

Passing to the limit $\Delta z \to 0$ we get: $\frac{dw}{dz} = \frac{1}{e^w} \implies \frac{d}{dz}(\text{Ln } z) = \frac{1}{z}$.
\end{proof}
\end{proposition}

\begin{example} \label{example:17.3} Study the function $w = \text{Ln}(z^2 + 4)$.
\begin{soln}
    
    $w = \text{Ln}(z^2 + 4) = \text{Ln}(z + 2i)(z - 2i) = \text{Ln}(z + 2i) + \text{Ln}(z - 2i)$ (Prop 15.1 \ref{proposition:15.1} 1)) $\implies w$ has 2 branch points $z = 2i$ and $z = -2i$. Therefore, $w$ is analytic in the cut-plane where we cut along 2 rays that start at $z = \pm 2i$ (including $\pm 2i$):

    \includegraphics{img_graph4.pdf}
    
    The derivative of this function in the cut-plane is computed as a composition function using: By the Chain Rule (Prop 6.1 6)): $(\text{Ln}(z^2 + 4))' = \frac{2z}{z^2 + 4}$.
\end{soln}
\end{example}

Now we define the main branch of the function $\text{Ln } z$.
\begin{definition}\label{definition:15.2}
    The branch of $w = \text{Ln } z$ when $k = 0, -\pi < \varphi \leq \pi$ is called the main branch of the logarithm: $w_0 = \text{Ln } z = \ln r + i\varphi$.
\end{definition}

Having the logarithm, similarly to the way that we have defined the trigonometric and hyperbolic functions using the exponential function, we define the inverse of the trigonometric and hyperbolic functions as follows:
\begin{definition}\label{definition:15.3}
    \begin{enumerate}
        \item $\arcsin z$ is $w \in \mathbb{C}$ such that $\sin w = z$.
        \item $\arccos z$ is $w \in \mathbb{C}$ such that $\cos w = z$.
        \item $\arctan z$ is $w \in \mathbb{C}$ such that $\tan w = z$.
    \end{enumerate}
    From 1): $z = \sin w$. By Euler's formula $z = \frac{e^{iw} - e^{-iw}}{2i}$.
\end{definition}

Eliminating $w$ we obtain: Denote $e^{iw} = t \implies z = \frac{t - \frac{1}{t}}{2i} \iff t^2 - 2izt - 1 = 0 \implies t = iz \pm \sqrt{1 - z^2}$ or $e^{iw} = iz \pm \sqrt{1 - z^2}$, when
 there are two branches of (each) square root. 
 Thus: $\arcsin z = w = \frac{1}{i} \text{Ln}(iz \pm \sqrt{1 - z^2}) = -i\text{Ln}(iz \pm \sqrt{1 - z^2})$.

In the same way:

\begin{exercise}
     $\arccos z = -i\text{Ln}(z \pm i\sqrt{1 - z^2})$
    $\arctan z = \frac{1}{2i} \text{Ln}\left(\frac{1+iz}{1-iz}\right)$.
\end{exercise}

\begin{example}\label{example:17.4} 
    Compute:
    \begin{enumerate}
        \item $\arcsin 2$
        \item $\arcsin i$
    \end{enumerate}
\begin{soln}
    \begin{enumerate}
        \item $\arcsin 2 = -i\text{Ln}(2i \pm \sqrt{1-4}) = -i\text{Ln}((2 \pm \sqrt{3})i) = -i(\ln(2 \pm \sqrt{3}) + i(\frac{\pi}{2} + 2\pi k)) = \frac{\pi}{2} + 2\pi k - i\ln(2 \pm \sqrt{3})$. If, for example, $k = 1$, $\arcsin 2$ attains two values: $\arcsin 2 \approx \frac{5}{2}\pi \pm i \cdot 1.32$.
        \item $\arcsin i = -i\text{Ln}(-1 \pm \sqrt{2})$. Let us compute each of the logarithms:
    \begin{align*}
        \text{Ln}(-1 + \sqrt{2}) &= \ln(\sqrt{2} - 1) + 2\pi k i \\
        \text{Ln}(-1 - \sqrt{2}) &= \ln(\sqrt{2} + 1) + (2k+1)\pi i
    \end{align*}
    $\implies \arcsin i = 2\pi k - i\ln(\sqrt{2} - 1), \quad k \in \mathbb{Z}$ ($k = 0, \pm 1, \pm 2, \dots$)
    $\arcsin i = \pi(2k+1) - i\ln(\sqrt{2} + 1), \quad k \in \mathbb{Z}$.
    \end{enumerate}
\end{soln}
\end{example}

\begin{remark}
    In a domains where the corresponding branches of the square root and of the logarithm are single-valued, the inverse trigonometric functions are analytic and
    \begin{align*} 
        & (\arcsin z)' = \frac{1}{\sqrt{1-z^2}} &\quad (\arccos z)' = -\frac{1}{\sqrt{1-z^2}} \\
        & (\arctan z)' = \frac{1}{1+z^2} 
    \end{align*}
\end{remark}


% \section*{Lectures 20+ 21.}
\lecture{20+21}{Singularities and other isolated ones...}

Now we introduce isolated singularities and discuss their classification. Then we define the residue of a function and relate calculation of residues of certain singularities to zeros of analytic functions.
\section{Isolated Singularities}

\subsection{Isolated singularities and definitions}

\begin{definition}\label{definition:16.1} 
    If $f$ is analytic on a punctured disc $D' = \{z \in \mathbb{C}: 0 < |z - z_0| < r\}$ but it is not analytic at $z_0$ (or not defined at $z_0$), then $z_0$ is called an isolated singularity of $f$.
\end{definition}

In other words, $z_0$ is an isolated singularity if there exists a neighborhood of $z_0$ such that $f$ is analytic at every point in this neighborhood except at $z_0$. 
If $z_0$ is an isolated singularity of $f$, then there exists an annulus $0 < |z - z_0| < r$ (sufficiently small) such that $f$ is analytic there 
and hence the Laurent series expansion of $f$ exists on this annulus.

By Thm~\ref{theorem:13.2} (Laurent series) Such $f$ has a Laurent series expansion valid on $D'$:
\[  f(z) = \sum_{n=1}^{\infty} b_n (z - z_0)^{-n} + \sum_{n=0}^{\infty} a_n (z - z_0)^n = \dots \frac{b_2}{(z-z_0)^2} + \frac{b_1}{z-z_0} + a_0 + \dots \]
The $\sum_{n=1}^{\infty} b_n (z - z_0)^{-n}$ part of the series is called the principal part of the Laurent series.

Note: $\sum_{n=0}^{\infty} a_n (z - z_0)^n \xrightarrow{z \to z_0} a_0$.

Now we can see that there are 3 options for the principal part in (\ref{equation:ln_z_formula}): there is no principal part; there is principal part and it has finitely many terms; or there is an infinite principal part. 
This observation leads to the following classification of isolated singularities.

\subsection{Three types of isolated singularities.}
I. Removable singularity: If $b_n = 0$ for all $n$, namely the principal part of the Laurent series is identically 0, then $z = z_0$ is called the removable singularity.

Namely, the Laurent series is an analytic function on the disc centered at $z_0$, including the point $z_0$. At every interior point of that disc the series converges to the function $f(z)$ and at $z_0$
($f(z)$ is not analytic at $z_0$) it converges to $a_0$. Define:
\[
 \tilde{f}(z) = \begin{cases}
f(z) & z \neq z_0 \\
a_0 & z = z_0
\end{cases}
\]
Then, $\tilde{f}(z)$ is the value of $\sum_{n=0}^{\infty} a_n (z - z_0)^n$ on $|z - z_0| < r$ (including $z = z_0$) and $\tilde{f}(z)$ is analytic at the point $z = z_0$.

II Pole: If the principal part of the Laurent series $\sum_{n=1}^{m} b_n (z - z_0)^{-n}$ has finite strictly positive number of non-zero terms, namely, 
if it has the form $\sum_{n=1}^{m} b_n (z - z_0)^{-n}$ with $b_m \neq 0$, then $z_0$ is called a pole of order $m$. A pole of order $m = 1$ is called a simple pole.

III Essential singularity: If $b_n \neq 0$ for infinitely many $n$, then $z_0$ is called an essential singularity.

\begin{definition}\label{definition:16.2}
     The coefficient $b_1$ of $\frac{1}{z - z_0}$ term is called the residue of the singularity at $z = z_0$, denoted by $\underset{z=z_0}{\text{Res}} f(z) = \text{Res}(f, z_0)$.
\end{definition}

First, let us study the case of removable singularity. The following proposition is a useful criterion to check for removable singularity.

\begin{proposition}\label{proposition:16.1}
    $z_0$ is a removable singularity of $f(z)$ if and only if $\lim_{z \to z_0} (z - z_0)f(z) = 0$.
\end{proposition}
\begin{proof}
 $\implies$ Suppose $z_0$ is removable and $g$ is analytic on the open disc centered at $z_0$ of radius $R: D = \{z \in \mathbb{C}: 0 \leq |z - z_0| < R\}$, such that $f = g$ for $z \neq z_0$. 
Then, since $g$ is continuous at $z_0$ (Prop 6.5) and \[ \lim_{z \to z_0} (z - z_0)f(z) = \lim_{z \to z_0} (z - z_0)g(z) = \lim_{z \to z_0} (z - z_0) \lim_{z \to z_0} g(z) = 0 \cdot g(z_0) = 0.\]

$\impliedby$ Suppose that $\lim_{z \to z_0} (z - z_0)f(z) = 0$ and $f$ is analytic on the punctured disc $D' = \{z \in \mathbb{C}: 0 < |z - z_0| < R \}$. Define:
\[ g(z) = \begin{cases}
(z - z_0)^2 f(z) & z \neq z_0 \\
0 & z = z_0
\end{cases}, \]
$g$ is analytic for $ z \neq z_0$ and it is differentiable at  $z_0$:
\[ g'(z_0) = \lim_{z \to z_0} \frac{g(z) - g(z_0)}{z - z_0} = \lim_{z \to z_0} \frac{(z - z_0)^2 f(z)}{z - z_0} = \lim_{z \to z_0} (z - z_0)f(z) = 0 \]
$\implies g$ is analytic on $D$ with $g(z_0) = g'(z_0) = 0 \implies g$ has power series expansion $g(z) = \sum_{n=0}^{\infty} a_n (z - z_0)^n$ with $a_0 = a_1 = 0 \implies$ we can factor $(z - z_0)^2$ 
from the series, so
\[ g(z) = (z - z_0)^2 \sum_{n=0}^{\infty} a_{n+2} (z - z_0)^n = (z - z_0)^2 \varphi(z).\] 
Therefore, for $z \neq z_0$ we have: $f(z) = \sum_{n=0}^{\infty} a_{n+2} (z - z_0)^n$ and this series defines an analytic function on $D$.
\end{proof}

The other types of singularities can also be recognised by the behavior of $f$ near it.
    \begin{enumerate}
        \item If $z_0$ is a pole, then $\lim_{z \to z_0} f(z) = \infty$.

If $z_0$ is a pole of order $m$, then using the Laurent series expansion we get
\begin{align*}
    f(z) &= b_m (z - z_0)^{-m} + \dots + b_1 (z - z_0)^{-1} + a_0 + a_1(z - z_0) + \dots \\ 
    & = \frac{1}{(z - z_0)^m} [b_m + b_{m-1}(z - z_0) + \dots + b_1(z - z_0)^{m-1} + a_0(z - z_0)^m + \dots] 
\end{align*}
Therefore, since $\lim_{z \to z_0} \frac{1}{(z - z_0)^m} = \infty$ and $\lim_{z \to z_0} [b_m + \dots] = b_m$, we conclude that $\lim_{z \to z_0} f(z) = \infty$.

        \item If $z_0$ is an essential singularity, then $\lim_{z \to z_0} f(z)$ does not exist - highly non-obvious!
    \end{enumerate}

In fact, more is true in the case of an essential singularity:

\subsection{Picard's theorem}
\begin{theorem}[The Great Picard Theorem]\label{theorem:great.picard}
    If an analytic function $f$ has an essential singularity at $z_0$, then on any punctured neighborhood of $z_0$, $f$ takes on all possible 
    values in $\mathbb{C}$ infinitely often with at most one exception.
\end{theorem}

In other words, if $z_0$ is an essential singularity of the function $f$, then the equation $f(z) = c$ for any (finite) $c$ except possibly for one value, has infinitely many solutions that tend to $z_0$.

Similarly to Prop (16.1) (\ref{proposition:16.1}), we have the following proposition which provides a method for recognising the order of a pole and computing the corresponding residue.

\begin{proposition}\label{proposition:16.2} 
    $z = z_0$ is a pole of order $m > 0$ of $f \iff f$ can be expressed as $f(z) = \frac{\varphi(z)}{(z - z_0)^m}$ where $\varphi(z)$ is analytic 
    at a neighborhood of $z_0$ (and at $z_0$) and $\varphi(z_0) \neq 0$. Moreover, $\text{Res}(f, z_0) = \frac{\varphi^{(m-1)}(z_0)}{(m-1)!}$.
\end{proposition}
\begin{proof}
$\implies$ If $z_0$ is a pole of order $m$, then as we have seen 
\[ f(z) = \frac{1}{(z - z_0)^m} [b_m + b_{m-1}(z - z_0) + \dots + b_1(z - z_0)^{m-1} + a_0(z - z_0)^m + \dots] \]

Denote: \[ \varphi(z) = b_m + b_{m-1}(z - z_0) + \dots + b_1(z - z_0)^{m-1} + a_0(z - z_0)^m + \dots.\]
Then, $\varphi$ is analytic at a neighborhood of $z_0$ (this is its Taylor series expansion!) including $z_0$ and $\varphi(z_0) = b_m \neq 0$.
$\impliedby$ Assume $f(z) = \frac{\varphi(z)}{(z - z_0)^m}$, where $\varphi$ is analytic in a neighborhood of $z_0$ (including $z_0$) and $\varphi(z_0) \neq 0$. Since $\varphi$ is analytic 
there exists Taylor series expansion in some neighborhood of $z_0$ given by
\[ \varphi(z) = \varphi(z_0) + \frac{\varphi'(z_0)}{1!}(z - z_0) + \dots + \frac{\varphi^{(m-1)}(z_0)}{(m-1)!}(z - z_0)^{m-1} + \dots\]
$\implies f$ can be written in terms of series expansion as 
\[ f(z) = \frac{\varphi(z_0)}{(z - z_0)^m} + \frac{\varphi'(z_0)}{(z - z_0)^{m-1}} + \dots + \frac{\varphi^{(m-1)}(z_0)}{(m-1)!} \frac{1}{z - z_0} + \frac{\varphi^{(m)}(z_0)}{m!} + \dots \]
This is a Laurent series with finitely many terms in its principal part. Since $\varphi(z_0) \neq 0$, $f$ has a pole of order $m$ at $z_0$. The residue is the coefficient of $\frac{1}{z - z_0}$, 
namely $\frac{\varphi^{(m-1)}(z_0)}{(m-1)!}$.
\end{proof}

Let us see examples of these singularities.

\begin{example}\label{example:18.1} Let $f(z) = \frac{1}{(z-1)(z-2)}$. Then $f$ has isolated singularities at $z = 1$ and $z = 2$ \{f is a rational function $\implies$ analytic everywhere except at the 
    zeros of the denominator\}.

\end{example}
\checkthis{} The Laurent series about the point $z_0 = 1$ is given by \[ f(z) = -\frac{1}{z-1} + \sum_{n=0}^{\infty} (-1)^n (z-1)^n.\]

The Laurent series at $z_0 = 2$: Observe that
\begin{align*} f(z) = \frac{1}{(z-2)(z-1)} & = \frac{1}{z-2} \cdot \frac{1}{1 + (z-2)} = \frac{1}{z-2} \sum_{n=0}^{\infty} (-1)^n (z-2)^n \\ 
    & = \frac{1}{z-2} + \sum_{n=0}^{\infty} (-1)^{n+1} (z-2)^{n-1}.
\end{align*}
The series is valid on a punctured disc centered at $z_0 = 2$ of radius 1: $D' = \{z \in \mathbb{C}: 0 < |z-2| < 1\}$.
\begin{claim}
    Claim: $z = 1, z = 2$ are simple poles.
\end{claim}
\begin{proof}
 $z = 2$: Rewrite $f(z)$ in the form $f(z) = \frac{\varphi(z)}{z-1}$, where $\varphi(z) = \frac{1}{z-1}$. Then $\varphi(z)$ is analytic in a 
 neighborhood of $z = 2$ (including $z = 2$) and $\varphi(2) = 1 \neq 0$, thus, by Prop~\ref{proposition:16.2} $z = 2$ is a simple pole.

$z = 1$: In the same way: $f(z) = \frac{\tilde{\varphi}(z)}{z-1}$, $\tilde{\varphi}(z) = \frac{1}{z-2}$, $\tilde{\varphi}$ is analytic in a 
neighborhood of $z = 1$ (including $z = 1$), and $\tilde{\varphi}(1) = -1 \neq 0$. Thus, by Prop~\ref{proposition:16.2} $z = 1$ is a simple pole.
\end{proof}
By Prop~\ref{proposition:16.2}: $\text{Res}(f, 2) = \frac{\varphi(2)}{0!} = 1; \text{Res}(f, 1) = \frac{\tilde{\varphi}(1)}{0!} = -1$.

\begin{example}\label{example:18.2} Prove:
    \begin{enumerate}
        \item $z = 3$ is a removable singularity of $f(z) = \frac{z^2 - 9}{z - 3}$.
        \item $z = 0$ is a removable singularity of $f(z) = \frac{\sin z}{z}$.
    \end{enumerate}

\end{example}

\begin{soln}
    \begin{enumerate}
        \item Since $\lim_{z \to 3} \frac{(z-3)(z+3)}{z-3} = 6$ and $\lim_{z \to 3} (z-3) \cdot \frac{(z-3)(z+3)}{z-3} = 0$ by Prop~\ref{proposition:16.1}$z = 3$ is a removable singularity. 
        Define: $\tilde{f}(z) = \begin{cases} \frac{z^2 - 9}{z-3} & z \neq 3 \\ 6 & z = 3 \end{cases} \implies \tilde{f}(z)$ is entire (analytic on $\mathbb{C}$).
        \item We prove this in two ways:
        \begin{enumerate}[label=\roman*)]
            \item The Laurent series for $f$ in $0 < |z| < R$ is $f(z) = \frac{1}{z}(z - \frac{z^3}{3!} + \frac{z^5}{5!} - \dots) = 1 - \frac{z^2}{3!} + \frac{z^4}{5!} - \dots$ $\implies$ the 
            principal part does not exist $\implies z = 0$ is a removable singularity.
            \item $\lim_{z \to 0} z \cdot \frac{\sin z}{z} = \lim_{z \to 0} \sin z = 0 \implies$ by Prop~\ref{proposition:16.1} $z = 0$ is a removable singularity.
        \end{enumerate}
    \end{enumerate}
\end{soln}

\begin{example}\label{example:18.3} Classify the singularities of $f(z) = e^{\frac{1}{z}}$ and compute the residue of $f$ at $z = 0$.
\begin{soln}
    
    Since $f(z)$ is analytic everywhere except for $z = 0$, $f$ has an isolated singularity at $z = 0$. To determine the nature of this singularity, we write the Laurent series expansion: 
    \[ e^{\frac{1}{z}} = e^{z^{-1}} = \sum_{n=0}^{\infty} \frac{1}{n!} z^{-n} = 1 + \frac{1}{z} + \frac{1}{2!} \cdot \frac{1}{z^2} + \frac{1}{3!} \cdot \frac{1}{z^3} + \dots\]
    $\implies$ the principal part $\sum_{n=1}^{\infty} b_n (z - z_0)^{-n}$ has infinitely many non-zero terms $\implies z = 0$ is an essential singularity. 
    By Picard's Great Theorem, $f$ takes every value of $\mathbb{C} \setminus \{0\}$ infinitely often. The coefficient of the term $\frac{1}{z}$ in the Laurent series expansion is 1, therefore $\text{Res}(e^{\frac{1}{z}}, 0) = b_1 = 1$.
\end{soln}
\end{example}
